\documentclass[11pt]{article}
\usepackage{geometry}
\usepackage{hyperref}
\geometry{margin=1in}

\title{PhD Proposal: Differential Privacy in Federated Learning}
\author{Sarah Kim \\ Advisor: Prof. Roberts}

\begin{document}
\maketitle

\section{Problem Statement}
Federated learning trains models across devices without centralising raw data. Differential privacy adds formal guarantees but introduces accuracy-utility trade-offs.

\section{Research Questions}
\begin{enumerate}
  \item How can the privacy budget be allocated adaptively across training rounds?
  \item What is the optimal combination of clipping and noise for gradient aggregation?
  \item Can we achieve privacy-preserving personalisation?
\end{enumerate}

\section{Proposed Work}

\subsection{Year 1}
Analyse existing adaptive budget allocation schemes and derive lower bounds.

\subsection{Year 2}
Develop and empirically validate a new scheme on benchmark datasets (FedMNIST, FedCelebA).

\subsection{Year 3}
Extend to personalised federated learning and write dissertation.

\section{Expected Contributions}
\begin{itemize}
  \item Theoretical: tighter lower bounds on adaptive privacy allocation.
  \item Empirical: reference implementation and benchmark suite.
  \item Applied: deployment in one or two industry collaborations.
\end{itemize}

\section{Timeline}
First results paper: Q4 Year 1. Benchmark release: Q2 Year 2. Personalisation paper: Q1 Year 3. Thesis draft: Q4 Year 3.

\section{Related Work}
Prior work has examined static privacy budgets~\cite{dp-fl-2020} and gradient clipping~\cite{clipping-2021}. Adaptive strategies remain under-explored.

\end{document}
